\section{Билет 31. Звуковые волны. Область слышимости. Порог слышимости, порог болевого ощущения. Уровень громкости. Шкала уровней силы звука. Высота звука.}

\begin{definition}
\textbf{Звуковыми волнами} (или просто звуком) называются упругие волны с частотой в диапазоне от 16 Гц до 20 кГц, способные вызывать слуховое ощущение у человека при распространении в воздухе.
\begin{itemize}
    \item \textbf{Инфразвук}: $\nu < 16$ Гц --- не воспринимается ухом, но может ощущаться телом (вибрация);
    \item \textbf{Слышимый звук}: $16 \text{ Гц} \le \nu \le 20 \text{ кГц}$;
    \item \textbf{Ультразвук}: $\nu > 20$ кГц --- используется в медицине, дефектоскопии, эхолокации.
\end{itemize}
Источниками звука являются тела, колеблющиеся в среде с определённой частотой.
\end{definition}

\begin{definition}
\textbf{Область слышимости} --- диапазон интенсивностей и частот, в котором человеческое ухо способно воспринимать звук.
\begin{itemize}
    \item \textbf{Порог слышимости} $I_0$ --- минимальная интенсивность звука, способная вызвать слуховое ощущение. Принимается $I_0 = 10^{-12}$ Вт/м² для частоты 1 кГц.
    \item \textbf{Порог болевого ощущения} --- интенсивность, при которой звук вызывает не слуховое, а болевое ощущение. Соответствует $I_{\text{бол}} \approx 1$ Вт/м² ($L \approx 120$ дБ).
\end{itemize}
Чувствительность уха зависит от частоты: максимальна в диапазоне 1--5 кГц (речь), снижается на низких и высоких частотах.
\end{definition}

\begin{definition}
\textbf{Уровень громкости} $L$ --- логарифмическая мера интенсивности звука, основанная на законе Вебера--Фехнера (субъективное ощущение громкости растёт пропорционально логарифму интенсивности):
\begin{equation}
    L = 10 \lg \frac{I}{I_0} \quad \text{[дБ]}, \label{eq:loudness_level_31}
\end{equation}
где $I_0 = 10^{-12}$ Вт/м² --- интенсивность на пороге слышимости. Единица измерения --- децибел (дБ), $1 \text{ Б} = 10 \text{ дБ}$.
\end{definition}

\begin{property}[Шкала уровней силы звука]
\begin{center}
\begin{tabular}{|c|c|l|}
\hline
$L$, дБ & $I$, Вт/м² & Характеристика / источник \\
\hline
0 & $10^{-12}$ & Порог слышимости \\
10 & $10^{-11}$ & Тихий шелест листьев \\
20 & $10^{-10}$ & Шёпот (на расстоянии 1 м) \\
40--50 & $10^{-8}$--$10^{-7}$ & Обычный разговор \\
70--80 & $10^{-5}$--$10^{-4}$ & Громкий разговор, уличный шум \\
90--100 & $10^{-3}$--$10^{-2}$ & Грузовой автомобиль, рок-концерт \\
120 & $1$ & Болевой порог, отбойный молоток \\
140 & $10^2$ & Взлетающий реактивный самолёт \\
\hline
\end{tabular}
\end{center}
При $L > 160$ дБ возможен разрыв барабанных перепонок, при $L > 200$ дБ --- летальный исход.
\end{property}

\begin{remark}
\textbf{Высота звука} --- субъективная характеристика, зависящая в первую очередь от частоты колебаний: чем выше частота, тем выше воспринимаемый тон. Однако высота также зависит от интенсивности и тембра.
\par\noindent\textbf{Тембр} --- качественная характеристика звука, определяемая спектром (набором обертонов) и их относительными амплитудами. Именно тембр позволяет различать звуки одинаковой высоты и громкости от разных источников (например, скрипка и флейта).
\end{remark}
